% Lösungen Einheit 2 von 4 – Winkel berechnen (zusammenbau v0.1)
\einheitenkopf{Einheit 2 von 4 · Winkel berechnen}

%% TODO Lösung der Erklärzeile 19g) fehlt in der Bank
\erg{19}{a) Winkel: Umkehrtaste \quad b) $\mathrm{sin}\,\alpha = \frac{3}{8} = 0{,}375$, $\alpha = \mathrm{sin}^{-1}(0{,}375) \approx 22{,}0^\circ$ \quad c) $\mathrm{sin}\,\alpha = \frac{5}{9} = 0{,}5556$, $\alpha = \mathrm{sin}^{-1}(0{,}5556) \approx 33{,}7^\circ$ \quad d) $\mathrm{sin}\,\alpha = \frac{2}{5} = 0{,}4$, $\alpha = \mathrm{sin}^{-1}(0{,}4) \approx 23{,}6^\circ$ \quad e) $\mathrm{sin}\,\alpha = \frac{7}{10} = 0{,}7$, $\alpha = \mathrm{sin}^{-1}(0{,}7) \approx 44{,}4^\circ$ \quad f) $\mathrm{sin}\,\alpha = \frac{5}{8} = 0{,}625$, $\alpha = \mathrm{sin}^{-1}(0{,}625) \approx 38{,}7^\circ$ \quad g) \ldots \quad h) $\mathrm{cos}\,\alpha = \frac{5}{7}$, $\alpha = \mathrm{cos}^{-1}\left(\frac{5}{7}\right) \approx 44{,}4^\circ$ \quad i) $\mathrm{tan}\,\alpha = \frac{5}{8}$, $\alpha = \mathrm{tan}^{-1}\left(\frac{5}{8}\right) \approx 32{,}0^\circ$ \quad j) $b$ ist Gegenkathete, $a$ Hypotenuse: $\mathrm{sin}\,\beta = \frac{3}{10}$, $\beta \approx 17{,}5^\circ$ \quad k) $\alpha \approx 36{,}9^\circ$, $\beta = 90^\circ - 36{,}9^\circ = 53{,}1^\circ$; Kontrolle: $\mathrm{cos}\,\beta = \frac{3}{5}$ ergibt ebenfalls $\beta \approx 53{,}1^\circ$ \quad l) $\mathrm{sin}\,\alpha = \frac{2{,}3}{4{,}7}$, $\alpha \approx 29{,}3^\circ$ \quad m) $\mathrm{sin}\,\alpha = \frac{420}{1\,300}$, $\alpha \approx 18{,}8^\circ$ \quad n) Gegenkathete $3$, Ankathete $4$: $\mathrm{tan}\,\beta = \frac{3}{4}$, $\beta \approx 36{,}9^\circ$ \quad o) Teildreieck mit der Höhe: $\mathrm{sin}\,\alpha = \frac{4{,}5}{7{,}2}$, $\alpha \approx 38{,}7^\circ$ \quad p) $\mathrm{cos}\,\alpha = \frac{5{,}2}{8}$, $\alpha = \mathrm{cos}^{-1}\left(\frac{5{,}2}{8}\right) \approx 49{,}46^\circ$, also rund $49^\circ$ \quad q) $\mathrm{cos}\,\alpha = \frac{460}{610}$, $\alpha \approx 41{,}1^\circ$}
\erg{20}{a) Nein – der kürzeren Kathete liegt der kleinere spitze Winkel gegenüber; er muss kleiner als der andere sein, also unter der Hälfte des rechten Winkels.}
\erg{21}{a) $\mathrm{sin}^{-1}(0{,}42) \approx 24{,}8^\circ$}
\erg{22}{a) Fehler: Katheten vertauscht, das ist der andere Winkel. Richtig: $\mathrm{tan}\,\alpha = \frac{3}{8}$, $\alpha \approx 20{,}6^\circ$ \quad b) Der Sinus ist Gegenkathete durch Hypotenuse; die Hypotenuse ist die längste Seite, darum ist der Wert nie größer als 1 – zu 1,3 gibt es keinen Winkel. \quad c) $\mathrm{sin}\,\alpha = \frac{0{,}32}{5}$, $\alpha \approx 3{,}7^\circ$ – steiler als $3{,}4^\circ$, nicht zulässig.}
